\documentclass[12pt]{article}
\usepackage{amsmath,amssymb,amsfonts}
\usepackage[margin=1in]{geometry}

\begin{document}

\section*{Problem 5.31}

\textbf{Problem:} Suppose that we wish to find an eigenvector and an eigenvalue of a self-adjoint operator $A = A_0 + \epsilon A_1 + \ldots$. Imagine that we know the resolution of the identity, $E_\lambda(A_0)$, belonging to $A_0$ and assume that for a sufficiently small any eigenvalue of $A$ can be written as
\[
\lambda = \lambda_0 + \epsilon\lambda_1 + \epsilon^2\lambda_2 + \cdots
\]
and that the corresponding eigenvector can be written as
\[
\phi = \phi_0 + \epsilon\phi_1 + \epsilon^2\phi_2 + \cdots,
\]
where $A_0\phi_0^{(0)} = \lambda_0^{(0)}\phi_0^{(0)}$. Show in the manner of Section 4.11 that we can find an inhomogeneous equation for $\phi_1^{(0)}$, and that the solution to this equation which satisfies $(\phi_1^{(0)}, \phi_0^{(0)}) = 0$ is
\[
\phi_1^{(0)} = \int \frac{dE_\lambda A_1\phi_0^{(0)}}{\lambda_0^{(0)} - \lambda},
\]
where the integral excludes $\lambda = \lambda_0^{(0)}$.

(i) Write down and prove a result analogous to Eq.\ (1). (ii) For $\lambda_1$, how would the result generalize for $\lambda_2$?

\bigskip
\textbf{Solution:}

\subsection*{Setup: Perturbation theory using spectral resolution}

We have $A = A_0 + \epsilon A_1$ and seek solutions of $A\phi = \lambda\phi$.

Expanding to first order in $\epsilon$:
\[
(A_0 + \epsilon A_1)(\phi_0 + \epsilon\phi_1 + \cdots) = (\lambda_0 + \epsilon\lambda_1 + \cdots)(\phi_0 + \epsilon\phi_1 + \cdots).
\]

\textbf{Zeroth order:}
\[
A_0\phi_0 = \lambda_0\phi_0. \tag{0}
\]

\textbf{First order:}
\[
A_0\phi_1 + A_1\phi_0 = \lambda_0\phi_1 + \lambda_1\phi_0. \tag{1}
\]

Rearranging:
\[
(A_0 - \lambda_0)\phi_1 = (\lambda_1 - A_1)\phi_0. \tag{2}
\]

\subsection*{Finding $\lambda_1$}

Taking the inner product of (2) with $\phi_0$:
\[
(\phi_0, (A_0 - \lambda_0)\phi_1) = (\phi_0, (\lambda_1 - A_1)\phi_0).
\]

Since $A_0$ is self-adjoint, $(\phi_0, A_0\phi_1) = (A_0\phi_0, \phi_1) = \lambda_0(\phi_0, \phi_1)$, so the left side vanishes. Therefore:
\[
\boxed{\lambda_1 = (\phi_0, A_1\phi_0) = \langle\phi_0|A_1|\phi_0\rangle.}
\]

\subsection*{Finding $\phi_1$ using the spectral resolution}

From equation (2), we need to solve $(A_0 - \lambda_0)\phi_1 = (\lambda_1 - A_1)\phi_0$ with the constraint $(\phi_0, \phi_1) = 0$.

Using the spectral resolution of $A_0$:
\[
A_0 = \int_{-\infty}^{\infty}\lambda\,dE_\lambda.
\]

Formally inverting $(A_0 - \lambda_0)$ on the orthogonal complement of $\phi_0$:
\[
\phi_1 = -(A_0 - \lambda_0)^{-1}(A_1 - \lambda_1)\phi_0,
\]
where the inverse is taken on the subspace orthogonal to $\phi_0$. Using the spectral resolution:
\[
\phi_1 = -\int_{\lambda \neq \lambda_0}\frac{dE_\lambda(A_1 - \lambda_1)\phi_0}{\lambda - \lambda_0} = \int_{\lambda \neq \lambda_0}\frac{dE_\lambda\,A_1\phi_0}{\lambda_0 - \lambda},
\]
where we used $dE_\lambda\phi_0 = 0$ for $\lambda \neq \lambda_0$ (since $\phi_0$ is an eigenstate of $A_0$ at $\lambda_0$), so the $\lambda_1$ term drops out.

\[
\boxed{\phi_1 = \int_{\lambda\neq\lambda_0}\frac{dE_\lambda\,A_1\phi_0}{\lambda_0 - \lambda}.}
\]

This satisfies $(\phi_0, \phi_1) = 0$ by construction (the integral excludes $\lambda_0$).

\subsection*{Part (ii): Second-order energy $\lambda_2$}

From the second-order equation:
\[
A_0\phi_2 + A_1\phi_1 = \lambda_0\phi_2 + \lambda_1\phi_1 + \lambda_2\phi_0.
\]

Taking the inner product with $\phi_0$:
\[
\lambda_2 = (\phi_0, A_1\phi_1) = \int_{\lambda\neq\lambda_0}\frac{|dE_\lambda A_1\phi_0|^2}{\lambda_0 - \lambda}.
\]

In Dirac notation, this is the familiar second-order result:
\[
\boxed{\lambda_2 = \int_{\lambda\neq\lambda_0}\frac{\|dE_\lambda A_1\phi_0\|^2}{\lambda_0 - \lambda} = \sum_{n\neq 0}\frac{|\langle n|A_1|0\rangle|^2}{\lambda_0 - \lambda_n}.}
\]

\end{document}
